\section{Appendix}\label{sec:appendix}

\subsection{Setting A vs.\ Setting B (Pilot)}

In an early 30-question pilot we measured the KG agents under
\emph{Setting A}, in which a single LLM played both reasoner and
query-author roles, before fixing the query agent at
Qwen3-Coder-30B (Setting B). Setting A produced strongly asymmetric
results across LLM families that confounded the headline KG-vs-
context comparison:

\begin{itemize}
\item Qwen3-235B-Instruct on \backendname{typedb}: 83\% in
      Setting A, 77\% in Setting B (--6 pp; the larger general
      model authored TQL more fluently than the smaller coder).
\item Qwen3-235B-Instruct on \backendname{neo4j}: 60\% in
      Setting A, 83\% in Setting B (+23 pp; Cypher is more
      represented in Qwen-Coder's training).
\item HCX on \backendname{typedb}: 60\% in Setting A, 63\% in
      Setting B (+3 pp).
\end{itemize}

The asymmetry confirms that single-LLM measurements conflate KG-
backend effects with query-language familiarity. Setting B isolates
the chat-model variable; all main-paper numbers are Setting B.

\subsection{Surface-Form Sample Excerpts}

For reference, the same fact —"a Grade-1 employee in the position
\emph{부서장(가)} (Department Head Type-1) has an annual
position-pay of 18{,}192{,}000 KRW"— is rendered in the four
in-context tiers as follows.

\paragraph{C-TQL.}
\begin{verbatim}
$g1 isa 직급, has 직급명 "1급", has 직급서열 1;
$p_bsg_a isa 직위, has 직위명 "부서장(가)";
$pp_P01_1급 isa 직책급기준, has 직책급액 18192000;
(해당직급:$g1, 해당직위:$p_bsg_a, 적용기준:$pp_P01_1급) isa 직책급결정;
\end{verbatim}

\paragraph{C-Cypher.}
\begin{verbatim}
CREATE (g1:JobGrade {name: "1급", order: 1});
CREATE (p_bsg_a:Position {code: "P01", name: "부서장(가)"});
CREATE (pp_P01_1급:PositionPay {annual_amount: 18192000});
CREATE (g1)-[:HAS_POSITION_PAY]->(pp_P01_1급);
CREATE (p_bsg_a)-[:HAS_POSITION_PAY]->(pp_P01_1급);
\end{verbatim}

\paragraph{C-NL.}
\begin{quote}
부서장(가) 직위에 있는 1급 직원의 연간 직책급은 18,192,000원이다. \\[2pt]
\textit{[gloss: A Grade-1 employee in the position
부서장(가) (Department Head Type-1) has an annual position-pay of
18{,}192{,}000 KRW.]}
\end{quote}

\paragraph{C-JSONLD.}
\begin{verbatim}
{
  "직책급": [
    {"직위코드": "P01", "직위명": "부서장(가)",
     "직급": "1급", "연간직책급": 18192000}
  ]
}
\end{verbatim}

\subsection{Cypher Syntax Guide (Section 6.4)}

The 453-character Cypher primer prepended to C-Cypher+guide for the
Llama-4-Scout ablation:

\begin{quote}\small
\textit{[The Cypher dump represents a Neo4j property graph. Read
\texttt{CREATE (var:Label \{prop: value, ...\})} as defining a node
of the given label with properties; read
\texttt{CREATE (a)-[:REL\_TYPE \{prop: value\}]->(b)} as a directed
binary relationship from \texttt{a} to \texttt{b} with optional
edge properties. Same variables refer to the same node across
statements. Multi-arity facts (e.g., position pay with both
position and grade) are flattened by introducing intermediate
nodes connected to all participants.]}
\end{quote}
